\section{数据处理与描述性分析}

\subsection{数据来源与清洗口径}

本文使用附件中的布伦特原油期货主力合约价格数据，字段包括交易日期、开盘价、最高价、最低价、收盘价和前收盘价。原始数据共 2237 条，覆盖 2017 年 9 月 1 日至 2026 年 5 月 5 日。清洗工作包括日期解析、数值字段标准化、收益率计算、滚动波动率计算和冲突窗口截取。第一条记录的前收盘价缺失属于正常时间序列边界，冲突窗口内收盘价不存在缺失。

本文采用 \texttt{close\_price} 作为短期拟合和误差评价目标；\texttt{open\_price}、\texttt{high\_price}、\texttt{low\_price} 与 \texttt{pre\_close\_price} 用于日内振幅、收益率、滚动波动率和冲击峰值描述。题面极端峰值与附件 CSV 的收盘价口径不同，因此题面峰值只作为冲击背景，不参与 RMSE、MAE、MAPE 等误差计算。

为统一数据口径，本文将数据分为“拟合目标”“长期约束”“风险校验”和“候选市场变量”四类。EIA、JODI、OPEC、GPR、OVX 和期限结构资料用于数量级约束、风险校验或后续扩展，不构造新的真实油价。\cref{tab:external-data-role} 给出各类数据在模型中的位置。

\begin{table}[!htbp]
\centering
\caption{外部数据来源、模型作用与使用边界}
\label{tab:external-data-role}
\small
\begin{tabularx}{\linewidth}{L{2.8cm}L{3.0cm}L{4.0cm}Y}
\toprule
来源 & 数据对象 & 在模型中的作用 & 使用边界 \\
\midrule
赛题附件 CSV & 2017--2026 年布伦特原油期货主力合约价格 & 短期冲击模型拟合目标、基准模型对比、历史滚动窗口检验 & 短期 RMSE、MAE、MAPE 均以附件收盘价计算 \\
EIA 周度石油数据 & 美国 SPR 库存、商业库存、原油产量 & 中长期油价调节模型中的库存缓冲、SPR 释放强度和非海湾供给恢复数量级校验 & 美国口径，不等同全球总量；不参与短期模型拟合 \\
JODI Oil World Database & 多国月度产量、库存、进出口、需求、炼厂产出 & 多国上报口径下的供需与库存约束，辅助校验中长期参数方向 & 有国家覆盖差异和发布滞后，只作同口径数量级校验 \\
OPEC MOMR & 全球供需平衡表、需求增量、非 DoC 供给增量、DoC 原油需求差额 & 约束中长期油价调节模型中的需求背景、供给恢复和尾部压力 & 正常基线预测，不替代霍尔木兹封锁冲击参数 \\
Caldara-Iacoviello GPR & 月度地缘政治风险指数 & 外部校验地缘风险溢价，并进行滞后相关检验 & 指数由新闻构造，避免同月解释同月价格；不作为短期同步拟合输入 \\
FRED/Cboe OVX & 原油 ETF 隐含波动率 & 约束长期状态转移压力、不确定性强度和尾部风险 & 反映期权市场波动预期，不直接预测日度价格方向 \\
ICE / EIA 报道 & Brent 多期限官方结算价候选来源、现货-近月期货价差公开证据 & 候选期限结构变量，用于未来约束近端短缺和库存压力 & 当前未获得完整可复现多期限曲线，暂不写入数值模型；后续只采用 ICE 官方多期限结算价口径 \\
\bottomrule
\end{tabularx}
\end{table}

\begin{table}[!htbp]
\centering
\caption{附件数据和冲突窗口关键统计}
\begin{tabular}{lr}
\toprule
指标 & 数值 \\
\midrule
原始记录数 & 2237 \\
全样本日期范围 & 2017-09-01 至 2026-05-05 \\
冲突窗口日期范围 & 2026-03-02 至 2026-05-05 \\
冲突窗口交易日数 & 46 \\
冲突窗口最低收盘价 & 78.07 \\
冲突窗口最高收盘价 & 114.06 \\
冲突窗口最高盘中价 & 119.50 \\
冲突窗口平均收盘价 & 99.72 \\
冲突窗口末日收盘价 & 110.31 \\
\bottomrule
\end{tabular}
\end{table}

\subsection{外部供需与风险约束}

为避免中长期参数只由冲突窗口拟合得到，本文引入 EIA、JODI 和 OPEC 三类官方数据，对 SPR、商业库存、供给恢复、需求背景和尾部压力进行数量级约束。三类数据的口径不同：EIA 为美国周度数据，JODI 为多国上报月度数据，OPEC 为全球供需基线预测。因此，它们只约束参数方向和合理范围，不替代附件价格序列。

\begin{table}[!htbp]
\centering
\caption{EIA 官方外生数据摘要}
\label{tab:eia-summary}
\small
\begin{tabularx}{\linewidth}{Y L{2.2cm} S[table-format=6.0] S[table-format=6.0] S[table-format=-5.1] S[table-format=-2.2]}
\toprule
指标 & 变化口径 & {冲突前} & {窗口末} & {窗口变化} & {日均变化} \\
\midrule
美国 SPR 库存 & 千桶；万桶 & 415441 & 392700 & -2274.1 & -36.10 \\
美国商业库存 & 千桶；万桶 & 439279 & 457182 & 1790.3 & 28.42 \\
美国原油产量 & 千桶/日；万桶/日 & 13696 & 13573 & -12.3 & -0.20 \\
\bottomrule
\end{tabularx}
\end{table}

EIA 数据显示，美国 SPR 释放强度明显低于赛题给定的国际协调释放上限，美国商业库存未在冲突窗口内快速耗尽，美国原油产量短期变化较小。这说明模型中的 SPR 参数应解释为多国协调缓冲能力；库存缓冲不能写成机械消耗项；非海湾供给恢复也难以单独抵消封锁缺口。

JODI Oil World Database 汇集多国官方上报的月度石油数据。本文抽取一次石油和二次石油中的产量、库存、进出口、油品需求与炼厂产出，并采用 2026 年 2 月已上报国家/地区作为同口径样本。由于 JODI 存在国家覆盖差异和发布滞后，本文将其表述为“多国上报口径”。

\begin{table}[!htbp]
\centering
\caption{JODI 多国口径外生约束摘要}
\label{tab:jodi-summary}
\small
\begin{tabularx}{\linewidth}{Y C{1.45cm} S[table-format=7.2] C{1.45cm} S[table-format=-5.2]}
\toprule
指标 & 月份 & {数值} & 单位 & {同口径较上年末变化} \\
\midrule
原油及凝析油产量 & 2026-02 & 51456.97 & KBD & 108.39 \\
原油及凝析油期末库存 & 2026-02 & 2203133.02 & KBBL & 20073.83 \\
原油及凝析油进口 & 2026-02 & 26746.71 & KBD & -963.69 \\
原油及凝析油出口 & 2026-02 & 22887.30 & KBD & -465.39 \\
油品总需求 & 2026-02 & 52784.43 & KBD & 423.81 \\
油品炼厂总产出 & 2026-02 & 54428.66 & KBD & 4316.20 \\
油品期末库存 & 2026-02 & 1837991.05 & KBBL & 49713.86 \\
\bottomrule
\end{tabularx}
\end{table}

JODI 数据提示，非海湾供给恢复和需求调整是多国系统的缓慢响应，库存缓冲也需要同时区分原油库存和成品油库存。由于 2026 年 2 月可得数据处于冲突发生前，它更适合作为中长期外推的先验约束。

OPEC Monthly Oil Market Report 给出全球正常情景下的年度和季度供需基线，适合约束中长期油价调节模型中的需求背景、非 DoC 供给恢复和 DoC 原油需求差额尾部压力。

\begin{table}[!htbp]
\centering
\caption{OPEC 全球供需平衡表关键约束}
\label{tab:opec-summary}
\small
\begin{tabularx}{\linewidth}{Y S[table-format=2.2] C{1.4cm} Y}
\toprule
指标 & {数值} & 单位 & 模型含义 \\
\midrule
2026 年世界石油需求增量 & 1.40 & mb/d & 长期需求背景仍增长，需求收缩不能设得过强 \\
2026 年非 DoC 液体供给增量 & 0.60 & mb/d & 外部供给有恢复能力，但远小于封锁冲击量级 \\
2026 年 DoC 原油需求差额增量 & 0.60 & mb/d & 对 DoC 原油的压力仍上升，支持保留尾部风险 \\
2026 年 DoC 原油需求差额峰值 & 43.70 & mb/d & 下半年需求差额更高，中长期油价调节模型需保留季节性压力 \\
峰值高于年均 & 0.70 & mb/d & 用于温和抬升尾部风险权重 \\
\bottomrule
\end{tabularx}
\end{table}

OPEC 表说明，在没有霍尔木兹极端封锁的正常基线下，2026 年全球需求仍有增长、非 DoC 供给恢复有限且下半年 DoC 原油需求差额更高。因此，中长期油价调节模型中的价格回落应解释为政策缓冲、供给恢复、需求弹性和风险溢价衰减共同作用后的条件路径，而非全球需求自然塌陷。

本文还接入 Caldara-Iacoviello 官方 GPR 月度指数和 Cboe/FRED 的 OVX 原油 ETF 隐含波动率。GPR 由新闻文本构造，因此不用于同月解释同月价格；OVX 反映期权市场对未来约 30 日油价波动的定价，冲突窗口均值为 87.07，最大值为 120.91，同长度历史窗口均值分位约 97.3\%，适合进入长期风险权重和不确定性强度约束。

\keypoint{本章的数据处理逻辑是：附件价格用于拟合和评价，EIA、JODI、OPEC 约束供需和库存数量级，GPR 校验地缘风险环境，OVX 约束市场隐含波动率。}

\begin{figure}[!htbp]
\centering
\begin{subfigure}[t]{0.78\linewidth}
  \centering
  \includegraphics[width=\linewidth]{布伦特原油长期价格走势.png}
  \caption{2017--2026 年长期走势}
\end{subfigure}
\vspace{0.35em}
\begin{subfigure}[t]{0.78\linewidth}
  \centering
  \includegraphics[width=\linewidth]{冲突窗口价格走势.png}
  \caption{2026 年冲突窗口路径}
\end{subfigure}
\caption{附件收盘价的长期背景与冲突窗口放大}
\end{figure}

\FloatBarrier

\subsection{波动率特征}

冲突窗口内，油价从 78.07 美元/桶快速上行至高位平台，并在后期出现再定价回落和反复震荡。这一现象说明，价格不只是由供应缺口单向推升，还受到政策缓冲、市场确认和预期修复的共同影响。滚动波动率图显示，冲突窗口附近波动明显放大，适合使用事件机制模型而不是普通平稳时间序列模型。

\begin{figure}[htbp]
\centering
\includegraphics[width=0.92\linewidth]{布伦特收益率与波动率.png}
\caption{收益率与滚动波动率}
\end{figure}
